\section*{Bab \ref{ch:b2:blood-circulation} --- Darah dan Sistem Peredaran}

\begin{solution}{exo:b2:blood-circulation:1}
\omterm{def:b2:blood-circulation:blood}{Plasma} \qty{55}{\%} (air, \qty{70}{g/L} protein, garam, zat gizi, gas,
hormon); sel \qty{45}{\%}. \omterm{def:b2:blood-circulation:blood}{Sel darah merah} $5\times
10^{12}$/L, cakram \qty{7}{\micro m}, pengangkutan oksigen; sel \omterm{def:b2:blood-circulation:blood}{darah}
putih $7\times 10^{9}$/L, \qtyrange{10}{15}{\micro m}, kekebalan;
\omterm{def:b2:blood-circulation:blood}{keping darah} $3\times 10^{11}$/L, serpihan \qty{2}{\micro m}, hemostasis.
\end{solution}

\begin{solution}{exo:b2:blood-circulation:2}
Ikan: jantung $\to$ insang $\to$ tubuh $\to$ jantung, satu rangkaian,
dengan tubuhnya menerima \omterm{def:b2:blood-circulation:blood}{darah} pada tekanan rendah yang tersisa setelah
insangnya. Mamalia: jantung kanan $\to$ paru-paru (\qty{25}{mmHg}) $\to$ jantung kiri $\to$ tubuh
(\qty{120}{mmHg}) $\to$ jantung kanan. Rangkaian gandanya memberi
tubuhnya tekanan tinggi tanpa memaparkan paru-parunya kepadanya, dan
membiarkan kedua rangkaiannya diatur secara terpisah.
\end{solution}

\begin{solution}{exo:b2:blood-circulation:3}
\omterm{def:b2:blood-circulation:vessels}{Arteri}: dinding tebal dengan lapisan elastik dan otot --- menghantarkan
\omterm{def:b2:blood-circulation:blood}{darah} pada tekanan tinggi dan meratakan denyutnya. \omterm{def:b2:blood-circulation:vessels}{Arteriol}: dinding
sebagian besar otot polos --- menetapkan hambatan dan membagikan
alirannya. \omterm{def:b2:blood-circulation:vessels}{Kapiler}: setebal satu sel endotelium --- pertukaran. \omterm{def:b2:blood-circulation:vessels}{Vena}:
dinding tipis yang dapat teregang dengan katup --- mengembalikan \omterm{def:b2:blood-circulation:blood}{darah}
pada tekanan rendah dan menyimpan sebagian besar \omterm{def:b2:blood-circulation:blood}{darahnya}.
\end{solution}

\begin{solution}{exo:b2:blood-circulation:4}
$70\times 72 = \qty{5}{L/min}$, \qty{7200}{L} tiap hari, berbanding
\qty{5}{L} \omterm{def:b2:blood-circulation:blood}{darah}: volumenya melewati jantungnya \num{1440} kali tiap
hari. Tidak ada organ yang dapat membuat atau menghabiskan tujuh ton
\omterm{def:b2:blood-circulation:blood}{darah} tiap hari; \omterm{def:b2:blood-circulation:blood}{darah} yang sama pasti kembali --- ia beredar.
\end{solution}

\begin{solution}{exo:b2:blood-circulation:5}
$\Delta P = 8\eta LQ/\pi r^{4} = 8\times 3\times 10^{-3}\times
0.4\times 8.3\times 10^{-5}/(\pi\times 2.44\times 10^{-8}) =
\qty{10}{Pa}$, \qty{0.08}{mmHg}: aortanya tidak berongkos apa pun;
tekanannya dibelanjakan di dalam \omterm{def:b2:blood-circulation:vessels}{arteriolnya}.
\end{solution}

\begin{solution}{exo:b2:blood-circulation:6}
$(15/12)^{4} = 2.4$: hambatannya naik 2.4 kali lipat, alirannya turun
menjadi \qty{41}{\%}. $(15/20)^{4} = 0.32$: hambatannya turun menjadi
sepertiga, alirannya naik 3.2 kali lipat.
\end{solution}

\begin{solution}{exo:b2:blood-circulation:7}
Aorta $83/3 = \qty{28}{cm/s}$; \omterm{def:b2:blood-circulation:vessels}{kapiler} $83/3000 =
\qty{0.028}{cm/s}$; lintasnya $0.1/0.028 = \qty{3.6}{s}$. Ketika
berolahraga keluarannya naik lima kali lipat dan luas \omterm{def:b2:blood-circulation:vessels}{kapiler} yang
terbuka barangkali tiga kali lipat, sehingga waktu lintasnya turun
menjadi kira-kira \qty{2}{s} --- masih cukup bagi oksigennya untuk pergi.
\end{solution}

\begin{solution}{exo:b2:blood-circulation:8}
Ujung \omterm{def:b2:blood-circulation:vessels}{arterinya} $35 - 25 = \qty{+10}{mmHg}$ (penyaringan); ujung \omterm{def:b2:blood-circulation:vessels}{venanya} $15
- 25 = \qty{-10}{mmHg}$ (penyerapan kembali). Dengan $P_c = 28$ pada ujung
\omterm{def:b2:blood-circulation:vessels}{venanya} nilai bersihnya adalah $+3$: cairannya tersaring sepanjang
seluruhnya dan tidak ada yang terserap kembali --- edema, yaitu pergelangan kaki yang bengkak pada gagal jantung.
\end{solution}

\begin{solution}{exo:b2:blood-circulation:9}
$40/66000 = \qty{0.61}{mol/m^3}$; $\pi = 0.61\times 8.314\times 310 =
\qty{1570}{Pa} = \qty{11.8}{mmHg}$; separuh albuminnya, \qty{5.9}{mmHg}.
Natrium menyeberangi dinding \omterm{def:b2:blood-circulation:vessels}{kapilernya} dengan bebas, sehingga
kepekatannya sama pada kedua sisinya dan ia tidak mengerahkan tekanan
osmosis melintasinya.
\end{solution}

\begin{solution}{exo:b2:blood-circulation:10}
$RC = \qty{2}{s}$: $120e^{-0.4} = \qty{80}{mmHg}$. $C = 1$: $RC = 1$,
$120e^{-0.8} = \qty{54}{mmHg}$ --- dan volume sekuncup yang sama
menaikkan tekanan sistoliknya lebih banyak di dalam aorta yang kaku.
\omterm{thm:b2:blood-circulation:windkessel}{Tekanan denyut} yang lebar berarti puncak yang lebih tinggi yang harus
dilawan jantungnya ketika menyembur, lebih banyak kerja dan lebih
banyak oksigen bagi keluaran yang sama.
\end{solution}

\begin{solution}{exo:b2:blood-circulation:11}
Beristirahat: $90/83 = \qty{1.08}{mmHg.s/mL}$; berolahraga: $105/417 =
\qty{0.25}{mmHg.s/mL}$, yaitu penurunan empat kali lipat. Karena $R \propto r^{-4}$,
$r$ telah naik sebesar $4^{1/4} = 1.41$: \omterm{def:b2:blood-circulation:vessels}{arteriolnya} telah melebar
sebesar \qty{41}{\%} secara rerata.
\end{solution}

\begin{solution}{exo:b2:blood-circulation:12}
Keselanjaran memancangkan kelajuan pada tiap arasnya lewat penampang
seluruhnya, dan pohonnya dibangun sedemikian rupa sehingga penampangnya
seribu kali penampang aortanya di dalam \omterm{def:b2:blood-circulation:vessels}{kapilernya} lalu kecil kembali
di dalam \omterm{def:b2:blood-circulation:vessels}{venanya}; \omterm{thm:b2:blood-circulation:poiseuille}{hukum Poiseuille} menaruh hambatannya di dalam
\omterm{def:b2:blood-circulation:vessels}{arteriolnya}, tempat otot dapat mengubahnya, dan meninggalkan pembuluh
yang lebar hampir bebas dari rugi. Sebuah \omterm{def:b2:blood-circulation:circuits}{peredaran terbuka} tidak
memiliki \omterm{def:b2:blood-circulation:vessels}{kapiler} untuk memperlambat \omterm{def:b2:blood-circulation:blood}{darahnya} di tempat yang tertentu,
tidak memiliki pembuluh untuk menetapkan sebuah hambatan, dan tidak
memiliki cara untuk mengarahkan alirannya ke satu organ --- ia hanya
dapat mengaduk.
\end{solution}

\begin{solution}{pb:b2:blood-circulation:1}
\textbf{1.} $70\times 72 = \qty{5.0}{L/min}$; \qty{7300}{L} tiap hari.
\textbf{2.} $7300/5 \approx 1450$ kali.
\textbf{3.} Seperempat dari \qty{57}{g} pada 72 denyut tiap menit: \qty{14}{g}
$\times 4320 = \qty{62}{kg}$ tiap jam --- kira-kira berat seorang lelaki.
\textbf{4.} Tidak ada makanan yang dapat memasok, dan tidak ada
jaringan yang dapat menghabiskan, berat seorang lelaki berupa \omterm{def:b2:blood-circulation:blood}{darah}
tiap jam; satu-satunya jalan keluarnya adalah bahwa \omterm{def:b2:blood-circulation:blood}{darah} yang sama
kembali ke jantungnya.
\textbf{5.} Seutas tali yang diikatkan pada lengannya menggembungkan
\omterm{def:b2:blood-circulation:vessels}{vena} di bawahnya, bukan di atasnya, sehingga \omterm{def:b2:blood-circulation:blood}{darah} di dalam \omterm{def:b2:blood-circulation:vessels}{vena}
bergerak menuju jantungnya; ketika \omterm{def:b2:blood-circulation:blood}{darah} di dalam sebuah \omterm{def:b2:blood-circulation:vessels}{vena} ditekan
menuju tangannya, ia berhenti pada katupnya dan tidak dapat didorong
kembali --- katupnya hanya mengizinkan aliran menuju jantungnya.
\textbf{6.} \omterm{def:b2:blood-circulation:vessels}{Kapiler} yang menyambungkan \omterm{def:b2:blood-circulation:vessels}{arteri} ke \omterm{def:b2:blood-circulation:vessels}{vena}; Malpighi
melihatnya di dalam paru-paru katak pada tahun 1661.
\textbf{7.} Luasnya $\pi\times 1.25^{2} = \qty{4.9}{cm^2}$; $83/4.9 =
\qty{17}{cm/s}$.
\textbf{8.} $8\times 3\times 10^{-3}\times 0.4\times 8.3\times
10^{-5}/(\pi\times 2.44\times 10^{-8}) = \qty{10}{Pa}$, \qty{0.08}{mmHg}.
\textbf{9.} $R_{1} = 8\times 3\times 10^{-3}\times 10^{-3}/(\pi\times
5.06\times 10^{-20}) = \qty{1.5e14}{Pa.s/m^3}$; terpasang sejajar,
$1.5\times 10^{14}/3\times 10^{6} = \qty{5.0e7}{Pa.s/m^3}$.
\textbf{10.} $8.3\times 10^{-5}\times 5.0\times 10^{7} = \qty{4200}{Pa}
= \qty{31}{mmHg}$.
\textbf{11.} \omterm{def:b2:blood-circulation:vessels}{Kapiler} yang terbuka $10^{10}$, masing-masing $\pi(3\times 10^{-6})^{2}
= \qty{2.8e-11}{m^2}$: \qty{0.28}{m^2} $= \qty{2800}{cm^2}$; $v =
83/2800 = \qty{0.03}{cm/s}$; lintasnya $0.1/0.03 = \qty{3.3}{s}$.
\textbf{12.} Hambatannya $\times(1/0.9)^{4} = 1.52$: \qty{47}{mmHg};
jantungnya harus menaikkan tekanan \omterm{def:b2:blood-circulation:vessels}{arterinya} sebesar \qty{16}{mmHg}
atau keluarannya turun sepertiga.
\textbf{13.} $+10$, $-10$, dan $0$ mmHg.
\textbf{14.} Penyaringannya $10\times 2 = \qty{20}{L}$ tiap hari;
penyerapan kembalinya $8.5\times 2 = \qty{17}{L}$.
\textbf{15.} \qty{3}{L} \omterm{prop:b2:blood-circulation:exchange}{getah bening} tiap hari.
\textbf{16.} $\pi_c = \qty{12.5}{mmHg}$: bersihnya $+22.5$ pada ujung
\omterm{def:b2:blood-circulation:vessels}{arterinya} dan $+2.5$ pada ujung \omterm{def:b2:blood-circulation:vessels}{venanya} --- penyaringan di mana-mana,
tanpa penyerapan kembali: edema.
\textbf{17.} Bersihnya $+10$ dan $+3$, reratanya $+6.5$: penyaringan
sepanjang seluruh \omterm{def:b2:blood-circulation:vessels}{kapilernya}, kira-kira \qty{26}{L} tiap hari; pembuluh
\omterm{prop:b2:blood-circulation:exchange}{getah beningnya} tidak sanggup membawanya dan cairannya menumpuk di
dalam jaringannya, kaki lebih dahulu.
\textbf{18.} $4\times 10^{10}\times 2\pi\times 3\times 10^{-6}\times
10^{-3} = \qty{750}{m^2}$; $t = x^{2}/2D = (2\times 10^{-5})^{2}/
(2\times 10^{-9}) = \qty{0.2}{s}$.
\textbf{19.} $(93 - 3)/83 = \qty{1.08}{mmHg.s/mL}$.
\textbf{20.} $RC = \qty{2.2}{s}$; $120e^{-0.8/2.2} = 120\times 0.69 =
\qty{83}{mmHg}$.
\textbf{21.} $RC = \qty{1.1}{s}$; $120e^{-0.73} = \qty{58}{mmHg}$:
\omterm{thm:b2:blood-circulation:windkessel}{tekanan denyutnya} melebar dari 37 menjadi \qty{62}{mmHg}.
\textbf{22.} $120e^{-0.3/2.2} = \qty{105}{mmHg}$: pada laju yang cepat
tekanannya nyaris tidak turun di antara denyutnya.
\textbf{23.} $C\Delta P = 2\times 20 = \qty{40}{mL}$ tersimpan;
\qty{30}{mL} lainnya mengalir keluar lewat \omterm{def:b2:blood-circulation:vessels}{arteriolnya} selama sistolnya sendiri.
\textbf{24.} Bilik yang berkerut memeras pembuluhnya sendiri sampai
tertutup pada sistolnya, sehingga otot jantungnya dialiri pada
diastolnya, oleh tekanan yang dipertahankan aortanya yang mengerut
kembali; aorta yang kaku, yang membiarkan tekanan diastoliknya runtuh,
membuat jantungnya kelaparan di antara denyutnya.
\textbf{25.} Keluarannya \qty{5}{L/min}; penurunan \omterm{def:b2:blood-circulation:vessels}{arteriolnya}
\qty{31}{mmHg}; saringan bersihnya \qty{3}{L} tiap hari menuju \omterm{prop:b2:blood-circulation:exchange}{getah
beningnya}; diastoliknya \qty{83}{mmHg} bagi $C = 2$ dan \qty{58}{mmHg}
bagi $C = 1$.
\end{solution}
